%!TEX root = ./../../thesis-trier.tex

\label{def:partial_function} 
\IN{partial function}
A \emph{partial function} $f \colon X \nrightarrow Y$
% is given by a unique set $D \subseteq X \nrightarrow Y$
%such that $f \colon X \to Y$ is a function
either associates with $x \in X$ a unique value, written $\applyfunc{x}{f}$, 
or is undefined for $x \in X$,
and $f(x) = y$ for $x \in X$ and $y \in Y$ means %denotes the fact
that $f(x)$ is defined and equals~$y$.
Correspondingly, $f(x) \ne y$ means 
that either~$f(x)$ is defined but it does not equal $y$ or that $f(x)$ is undefined.
% When writing $\applyfunc{x}{f} = y$ for $x \in X$
%and $y \in Y$ we mean that $\applyfunc{x}{f}$ is defined and equals $y$.
Setting $D = \{ x \in X \mid \applyfunc{x}{f} \text{ is defined } \}$,
we obtain a (total) function $f \colon D \to Y$. 
Every  function can be seen as a partial function defined everywhere.
Given $f \colon X  \nrightarrow Y$ and $g \colon Y  \nrightarrow Z$,
their \emph{composition} is $\concat{f}{g} \colon X \nrightarrow Z$
where
$\applyfunc{x}{(\concat{f}{g})} = z$ for $x \in X$ and $z \in Z$
if and only if $\applyfunc{x}{f}$ is defined with $\applyfunc{\applyfunc{x}{f}}{g} = z$,
%are defined and $\applyfunc{\applyfunc{x}{f}}{g} = z$,
and otherwise $\applyfunc{x}{(\concat{f}{g})}$ is undefined. Other notions also carry over from functions without difficulty,
for example the image of a subset $S \subseteq X$
by setting $\applyfunc{S}{f} = \{ \applyfunc{x}{f} \mid x \in X \text{ and } \applyfunc{x}{f} \text{ is defined } \}$.
\IN{composition!of partial functions}